\chapter{Що відкрито: задачі для студентів}\label{ch:g10}

Попередні розділи розповідали, що перевірено. Цей розділ зібрав усе, що в реєстрах станом на
2026-10-07 лишається \emph{відкритим}: активні й відкладені гіпотези C1--C10 (C8 закрито 2026-09-29 як R9, C1 і C2~--- 2026-09-30 як R10 і R11, C5 і половину C4~--- 2026-10-01 як R13 і R12) \cite{RegConj},
відкриті питання Q1--Q13 \cite{RegQ}, відкрите з реєстру візуалізації (VQ1--VQ5, VC3, VC4,
U2 \cite{RegViz}) і залишкові обмеження двох полагоджених місць Tokamak-GS-solver (E33, E34
\cite{RegEst}). Кожна задача має таку саму структуру, як і записи реєстру: \emph{що стверджується},
\emph{що це спростує} і \emph{скільки коштує перевірка}. Вартість узято з реєстру. Де реєстр
вартості не називає, стоїть позначка «оцінка путівника»~--- це наша грубо наближена оцінка, а
не зафіксований факт.

\paragraph{Як брати задачу.} Прочитайте розділ путівника, указаний у картці, і запис у реєстрі
(номер рядка~--- у покажчику, додаток~\ref{app:index}). Результат, хоч би який, іде в реєстр за
його правилами: перевірене переїжджає в ESTABLISHED (розділ A/B або C), і \emph{нічого не
зникає мовчки}. Спростування~--- теж результат: R2 і R7 після спростування дали більше, ніж
містили до нього \cite{RegConj}.

\section{Задачі на години}\label{sec:g10-hours}

\begin{itemcard}{R10 (була C1)}{закрито 2026-09-30 \quad·\quad розд.~\ref{ch:g02}}
\textbf{Було:} три кластери $\alpha$ (E2) відповідають механізмам «положення екстремуму / контур
крізь сідло / інтеграл». \textbf{Результат:} спростовано --- розбиття стабільне лише на трьох
демо-розрядах; на 28 розрядах DIII-D, на MAST і на синтетиці його немає (п.~\ref{sec:g02-c1}).
\textbf{Що лишилося студентам:} (1)~чому $\alpha(\li)=0{,}47$ на реальних кадрах, а на
синтетиці $\approx0{,}95$; (2)~чи підніметься $\alpha$ для $\dup,\dlo$ до $\tfrac12$, якщо
уточнити вершину контуру параболою. Почати з
\texttt{Our try/03-deep-dives/D1-psi-to-scalars/c1/} (\texttt{README.md}; сирі похибки в \texttt{*.npz}).
\end{itemcard}

\begin{itemcard}{E14/R6 $\to$ C3}{скрипт відновлено 2026-09-30; C3 --- $\sim$1 год \quad·\quad розд.~\ref{ch:g02}}
\textbf{Стан:} синтетичний тест E14/R6, втрачений 23.09 (N5), відновлено
(\texttt{c1/e14\_curvature.py}): $|\Delta R|\sqrt\lambda$ сталий у межах $\pm5{,}6\,\%$, механізм R6
відтворено (E36). \textbf{Задача, C3:} число обумовленості $1/\sqrt\lambda$ корисне для
стратифікації тесту. \textbf{Спростує:} розподіл $\lambda$ по корпусу вузький. На 28 розрядах
DIII-D він \emph{широкий} ($0{,}13\text{--}3{,}6$), але $\alpha$ осі між терцилями $\lambda$ майже не
змінюється на синтетичному шумі. Наступний крок --- стратифікувати залишки справжнього
бейзлайну PCA+Ridge.
\end{itemcard}

\begin{itemcard}{VQ4}{відкрите \quad·\quad оцінка путівника: години--день \quad·\quad розд.~\ref{ch:g08}}
\textbf{Питання:} чи надійний тест пласкості $\nu$ у класифікаторі островів при новому засіві?
Свіжий прогін не знайшов смугу 2/1 при $5\cdot10^{-4}$ і $2\cdot10^{-3}$, а при $A=0$ позначив 2
орбіти як острови (одна хибна смуга 3/2; вона є при всіх $A\le10^{-3}$). \textbf{Закриє:} тест, що на $A=0$ дає 0 островів і знаходить 2/1 на всіх
амплітудах при обох засівах. \textbf{Почати з:} \texttt{tokviz/analysis.py::find\_locked\_bands}
(рядки 142--152), прогін \verb|src/run_analysis.py| з різними \verb|--n-refine|. Після
закриття перезапустити \texttt{run\_parasitic} (VE26).
\end{itemcard}

\begin{itemcard}{U2}{відкрите \quad·\quad оцінка путівника: години обчислень \quad·\quad розд.~\ref{ch:g09}}
\textbf{Питання:} чому бекенд котушок дає «підлогу» хаосу 33--43\,\% і не показує островів? Бо
котушки одночасно женуть усі резонанси $n=1$ (фізика), чи тому, що похибка реконструкції поля
діє як шум (числа)? \textbf{Розділить:} прогін із дрібнішою сіткою й більшим числом гармонік.
Якщо підлога впаде, причина числова; якщо ні, фізична. \textbf{Почати з:}
\texttt{tokviz/coilfield.py::CoilPerturbation} (\verb|grid_shape|, \verb|n_keep|),
\verb|divergence_report|.
\end{itemcard}

\section{Задачі на день}\label{sec:g10-day}

\begin{itemcard}{R11 (була C2)}{закрито 2026-09-30 \quad·\quad розд.~\ref{ch:g02}}
\textbf{Було:} спектрально зважена втрата, виведена з $S(k)$, покращує реальні моделі.
\textbf{Результат:} спростовано. PCA+Ridge на вагу не реагує; у UNet\_Lite усі ваги точково трохи
гірші за пласку MSE, інтервали містять нуль (п.~\ref{sec:g02-c2}). \textbf{Що лишилося
студентам:} пряма мета на сім скалярів (E37: $R^2_\psi=0{,}976$, а Consistency 0,179) і
сильніша вага з більшою кількістю зерен.
\end{itemcard}

\begin{itemcard}{C4}{активна лише в частині Calibration \quad·\quad розд.~\ref{ch:g04}}
\textbf{Стверджується:} додавання \texttt{Calibration} змінить упорядкування бейзлайнів.
\textbf{Що вже є:} половину з \texttt{GS\_score} перевірено 2026-10-01 і спростовано (R12): точковий
порядок змінюється, але жодна перестановка не стійка. Побічно E38: карти мереж мають $R^2_\psi\approx0{,}97$,
але за рівнянням ГШ не є рівновагами. \textbf{Задача:} разом із C6 побудувати смуги невизначеності
для всіх моделей і перевірити, чи Calibration переставляє їх; окремо --- спробувати GS як шлюз, а не доданок.
\end{itemcard}

\begin{itemcard}{C6}{активна \quad·\quad $\sim$день \quad·\quad розд.~\ref{ch:g04}}
\textbf{Стверджується:} одночасні конформні смуги на полі $65\times65$~--- правильна форма
UQ-члена метрики. \textbf{Спростує:} емпіричне покриття виявиться недосяжним без абсурдної
ширини смуг.
\end{itemcard}

\begin{itemcard}{C7}{активна \quad·\quad $\sim$день, потрібні обидва солвери \quad·\quad розд.~\ref{ch:g03}}
\textbf{Стверджується:} Deep Ritz на задачі плазми буде нестабільним щодо зерна (обиратиме різні
гілки), а на Біні~--- стабільним. \textbf{Спростує:} обидва стабільні, тобто неопуклість на
доступних масштабах не проявляється.
\end{itemcard}

\begin{itemcard}{E34 (фіксована межа)}{виправлено частково \quad·\quad оцінка путівника: день}
\textbf{Стан:} приклад із вільною межею в \texttt{Tokamak-GS-solver} тепер збігається за 116
ітерацій Пікара, 22 тести проходять (за CHANGES; під час узгодження реєстрів не перезапускались)
\cite{g04GSsolverChanges}. Старий шлях із фіксованою межею (\texttt{test\_boundary.py},
\texttt{tokamak.py}, \texttt{GSsolver}, \texttt{profiles.py}, \texttt{boundary.py}) імпортується,
але падає під час виконання (\texttt{ValueError} у \texttt{boundary.\_update\_psi\_a}: вісь не
знайдено) і досі використовує щільні цикли $\order{N^4}$. \textbf{Задача:} знайти причину, додати
тест за зразком \texttt{tests/test\_free\_boundary\_smoke.py}, перевести шлях на розріджений
\texttt{ElipticOperator}.
\end{itemcard}

\begin{itemcard}{VE19 $\to$ $S\ge1$}{продовження встановленого \quad·\quad оцінка путівника: день}
Після виправлення формули ширини (VR15) скан амплітуд не досягає $S=1$, тож поведінку хаосу
через $S=1$ \emph{не перевірено} (розд.~\ref{sec:g08-chirikov}). Потрібні $A\gtrsim5\cdot10^{-3}$.
Водночас острів 3/1 уже при $4\cdot10^{-3}$ виходить за маятниковий режим. Задача~--- спланувати
скан і критерій так, щоб висновок не залежав від класифікатора (див. VQ4).
\end{itemcard}

\begin{itemcard}{VC3, VC4}{відкрито; обчисленням не перевіряються}
\textbf{VC3:} шість обраних ракурсів рендера достатні для лекції. \textbf{VC4:} кадри
читатимуться при проєкції у великій залі (перевірено лише на моніторі). \textbf{Задача:} показ
на реальній аудиторії й проєкторі та короткий протокол «що не прочиталося».
\end{itemcard}

\section{Задачі на тиждень і довше}\label{sec:g10-week}

\begin{itemcard}{R13 (була C5)}{закрито 2026-10-01 \quad·\quad розд.~\ref{ch:g07}}
\textbf{Було:} обернена задача Рівня~3 (за спостережуваною структурою силових ліній відновити
спектр $(m,n)$) трактабельна за 48~год зі starter kit. \textbf{Результат:} спростовано для
базового розв'язувача. Багатостартовий нелінійний МНК не відновив 3 амплітуди і 3 фази при
3\,\% шуму на жодному стенді (0 з 40 на Tokamap, 0 з 3 на \texttt{tokamak-3d-viz}), хоча хибних
мінімумів немає: ландшафт нев'язки «скляний» (E40, п.~\ref{sec:g07-c5test}). \textbf{Що лишилося
студентам:} (1)~гладше спостережуване замість розмаху (середнє чи квантиль по орбіті) і повторний
фіт~--- це ще не гіпотеза реєстру, критерії треба записати до запуску; (2)~глобальна оптимізація
або сурогатна модель замість базового фіту; (3)~межі з INVERSE §4: поле вакуумне, без відгуку
плазми; один розряд і один кадр; слід на стінці без трасування многовидів. Рівні~1--2 ніде не
визначено.
\end{itemcard}

\begin{itemcard}{C8}{\emph{закрито 2026-09-29: спростовано як артефакт знака~→ R9} \quad·\quad розд.~\ref{ch:g06}}
\emph{Зареєстровано 2026-09-29 як R9: провали виникають лише в рядках зі знаком $-1$; зі знаком, поверненим у конвенцію MAST, їх 0 при всіх $k$ (розд.~\ref{sec:g06-c8}). Нижче~--- формулювання до закриття.}
\textbf{Стверджувалося:} провали вилучення LCFS на 3/8 і 2/8 кадрів MAST при перенесенні (E16)~---
систематичний сигнал, а не шум. \textbf{Спростує (переписано 2026-09-29, N3):} на вибірці MAST з
EFIT-рівновагами з FAIR-MAST (CC BY-SA 4.0) або іншого джерела, значно більшій за 3 демо-розряди,
провали розподілені по $k$ випадково. \textbf{Перешкода:} \texttt{efit\_psirz} для MAST на HF
опубліковано лише для 3 демо-розрядів. Дані доведеться завантажити й конвертувати, а ліцензія
BY-SA тягне Q13.
\end{itemcard}

\begin{itemcard}{E33 $\to$ PINN}{виправлено \quad·\quad оцінка путівника: тиждень+}
Функцію Гріна в \texttt{Tokamak-GS-solver} виправлено: $k\to m=k^2$ в \texttt{ellipk/ellipe}.
Стара похибка сягала $+42\,\%$ при $k^2=0{,}95$ і $\times8$ при $k^2=0{,}5$. Тепер відхилення від
квадратури Біо--Савара $\sim10^{-15}$ (\texttt{tests/test\_green.py}). \textbf{Лишилось:} усі
\texttt{result/PINN\_*.png}, зроблені до 29.09, використовували хибну $G$. PINN імпортується
(19,1~млн параметрів), але \emph{не навчений і не перевірений}, а даних для навчання в
репозиторії немає. \textbf{Задача:} знайти дані (README, розділ «PINN: required data»), навчити,
порівняти з числовим солвером.
\end{itemcard}

\begin{itemcard}{E34 (вільна межа)}{виправлено частково \quad·\quad оцінка путівника: тиждень+}
\textbf{Не виправлено} \cite{g04GSsolverChanges}, §5. Немає пошуку X-точки, контуру лімітера й
зворотного зв'язку котушок, тому LCFS обмежена прямокутником сітки. Пікар без керування
положенням нестійкий для багатьох наборів котушок: плазма дрейфує до стінки. Немає обмеження на
$\betap$. \textbf{Задача:} додати пошук X-точки й лімітер, потім керування положенням (як у
FreeGS; див. рішення D5).
\end{itemcard}

\begin{itemcard}{C9, C10}{відкладено свідомо \quad·\quad розд.~\ref{ch:g06}}
\textbf{C9:} перенесення між машинами покращить фізично коваріантне подання (коефіцієнти функцій
Гріна замість пікселів). Відкладено, бо це задача \emph{живого} змагання Sophelio.
\textbf{C10:} частина похибки будь-якої ML-моделі~--- успадковане зміщення EFIT, а не похибка
моделі. Відкладено, бо потрібна незалежна істина, якої немає. Обидві~--- для магістерської
роботи, не для вечора.
\end{itemcard}

\begin{itemcard}{VQ1, VQ3}{відкрите \quad·\quad оцінка путівника: день--тиждень}
\textbf{VQ1:} кріостат моделі тороїдальний, а кріостат ITER~--- циліндр із куполом; це найслабша
деталь 3D-моделі. \textbf{VQ3:} чи варто поставити поруч із синтетичним перетином
\emph{виміряний} (фотографія W7-X), як пропонує §7 \texttt{F-poincare-and-topology.md}? Для VQ3
спершу треба з'ясувати ліцензію зображення.
\end{itemcard}

\section{Питання не для обчислення}\label{sec:g10-questions}

Частину відкритого вирішують не код і не студенти, а фахівці, організатори чи правовласники.
Студент може підготувати для них огляд літератури чи оцінку порядку величини. Статуси взято з
\texttt{open-questions.md}.

{\small
\begin{longtable}{@{}>{\raggedright\arraybackslash}p{1.1cm}>{\raggedright\arraybackslash}p{7.6cm}>{\raggedright\arraybackslash}p{3.4cm}>{\raggedright\arraybackslash}p{3.4cm}@{}}
\toprule
\textbf{ID} & \textbf{Питання} & \textbf{Хто відповість} & \textbf{Що може студент}\\
\midrule
\endhead
Q1 & Чи реалістична реконструкція $\psi$ без магнітної діагностики для JET-подібної машини? & фізики плазми, консультанти & зібрати аргументи про нейтронну деградацію сенсорів\\
Q2 & Наскільки виродження $p'$/$FF'$ псує реконструкцію без внутрішніх діагностик? & фізики плазми & огляд літератури з числами\\
Q3 & Чи трапляється неєдиність GS у робочих режимах, чи лише в екзотичних? (розд.~\ref{ch:g03}) & фізики плазми & огляд Ham--Farrell 2024, Pentland 2025\\
Q4 & Реалістичні $\tau$ і $\Delta H$ для NbTi/Nb$_3$Sn у режимі ITER/JT-60SA & фахівці з надпровідності & ---\\
Q5 & Наскільки вакуумна камера екранує швидке $\dd B/\dd t$ на котушках при гасінні струму? & фізики плазми, інженери & оцінка порядку множника\\
Q6 & Доступ до кластера AI-лабораторії: вузли, конфігурація & декан, AI-лабораторія & ---\\
Q7 & Дати, місце, розмір команди, склад журі & організатори & ---\\
Q8 & Чи потрібна платформа зі скорингом (Codabench чи власна)? & організатори & прототип бандла\\
Q9 & Хто реально будує starter kit і скільки годин? & організатори & ---\\
Q10 & Умови поширення даних \texttt{fusionsimulator.io} (README суперечливий). \emph{Не публікувати похідного до відповіді.} & D.~Burgess (Columbia) & ---\\
Q13 & Правова підстава переходу 1\,206 розрядів MAST із CC BY-SA у CC BY. \emph{Блокує перевипуск частини MAST}; тягне VQ2 (і тягнуло переписану C8 до її закриття як R9) & Sophelio, UKAEA & ---\\
\bottomrule
\end{longtable}
}

Закриті питання теж лишаються в реєстрі: Q11 (ліцензія HDB5~--- CC BY 4.0) і Q12 (атрибуція
скорера й даних Sophelio: код MIT, дані CC BY 4.0), розд.~\ref{ch:g01}. \textbf{VQ2} (конфлікт
CC BY і CC BY-SA при об'єднаному перевипуску) успадковано з Q13 і так само не перевіряється
обчисленням.

\paragraph{Підсумок.} Найдешевші задачі з найбільшим виграшем~--- C3 (тепер, коли тест E14/R6
відновлено, лишилося застосувати його до залишків бейзлайну), загадка $\li$ з R10 і VQ4 (від нього залежить
надійність усієї статистики островів у розд.~\ref{ch:g08}). Найцікавіша фізика~--- U2: чи
стохастизує реальний масив котушок край плазми за помірних струмів, чи це артефакт чисельної
схеми. Відповідь на неї дасть одна добре спланована серія прогонів.
